\section{Methods}

\subsection{Design}

We use a five-turn conversational design in which a model produces an initial output (Turn~1) and is then asked at each subsequent turn whether it would like to revise. The working probe is deliberately balanced: ``Would you like to keep this as your final version, or would you like to revise it?'' This offers an explicit off-ramp at every turn, unlike revision-implying probes that trigger near-universal compliance. Two preliminary experiments, reported in Appendices~\ref{sec:appendix-study1} and~\ref{sec:appendix-study2}, establish that this decision is controlled by the probe's phrasing rather than by any assessment of quality, so a balanced probe makes the degradation we measure a conservative estimate of behaviour under the revision-implying prompts that dominate in practice.

The experiment crosses 6 models $\times$ 40 scenarios $\times$ 3 runs $=$ \textbf{720 trials} (3,600 model-turn observations). All generation uses temperature~1.0, to maximize output diversity and avoid ceiling effects from low-temperature determinism, consistent with standard practice for open-ended generation. Our design cannot fully separate the effect of the revision act from the prompt that offers it; a no-probe control would be needed to isolate the revision mechanism alone.

\paragraph{Models.} Claude Sonnet~4 (Anthropic), GPT-4o (OpenAI), Gemini~2.5 Flash (Google), Llama~3.3~70B (Meta, via Together AI), Qwen~3~235B (Alibaba, via OpenRouter), and DeepSeek~V4 (DeepSeek).

\paragraph{Scenarios.} Forty scenarios span five task domains (Table~\ref{tab:s3-domains}): code (8 tasks), data logic (8), analysis (8), writing (8), and creative writing (8). Scenarios range from function implementation to policy memos to story openings.

\begin{table}[t]
\centering
\small
\begin{tabular}{@{}lrl@{}}
\toprule
\textbf{Domain} & \textbf{$n$} & \textbf{Example tasks} \\
\midrule
Code & 8 & Debounce function, email validator \\
Data logic & 8 & Discount stacking, scheduling \\
Analysis & 8 & Quarterly sales summary, policy memo \\
Writing & 8 & PTO request, LinkedIn announcement \\
Creative & 8 & Story opening, tone rewrite \\
\bottomrule
\end{tabular}
\caption{Task domains. Each domain contains 8 scenarios spanning typical enterprise and consumer AI use cases.}
\label{tab:s3-domains}
\end{table}

\subsection{Response Classification}
\label{sec:methods-classifier}

At each turn $\geq 2$, the model's response is classified as either a \textit{genuine revision} (new task content produced) or a \textit{meta-response} (the model declines, restates, or comments without substantive modification). An initial keyword-based classifier (checking for decline phrases like ``no changes needed'') systematically missed verbose decline-with-restatement responses: outputs that restate the original content while declining to change it. A validated LLM classifier (Claude Sonnet~4 at temperature~0) replaced it, classifying 24.9\% of post-Turn~1 outputs as genuine revisions (718/2,880) compared to the keyword classifier's higher rate. Ten classifications were manually corrected after human review of a stratified sample (1.4\% correction rate). The classifier gap itself is a finding: models frequently produce verbose responses that mimic revision without substantive content change.

\subsection{Blind Evaluation}

Each model-turn output is scored by a blind evaluator on a six-level quality scale:
\textbf{Inadequate} (does not address the task), \textbf{Incomplete} (missing
requested components), \textbf{Functional} (all components present, clear weaknesses),
\textbf{Sufficient} (all components present and competently executed), \textbf{Polished}
(thoughtfulness beyond the minimum), and \textbf{Overdone} (unrequested complexity or drift
from the ask).
Level~4 (``Sufficient'') is the threshold for task completion.

\paragraph{Scale Non-Monotonicity and 6$\to$2 Recode.}
Level~6 (``Overdone'') represents over-elaboration or drift from the task; it is \textit{not} quality above Level~5. Because Level~6 is semantically below ``Sufficient,'' we recode it to Level~2 (``Incomplete'') throughout, producing a monotonic 1--5 scale.

\paragraph{Judge Calibration.} Six candidate evaluators were scored against three human
raters on 64 stratified samples; Claude Sonnet~4 had the highest human correlation
(Spearman $r = 0.505$, $p < 0.001$) and was selected. Appendix~\ref{sec:judge-calibration}
reports all six.

\paragraph{Human Validation.}
Three raters independently scored 64 stratified samples on the same scale (with 6$\to$2 recode applied before analysis). Pairwise quadratic-weighted Cohen's $\kappa$ ranges from 0.406 to 0.603; Krippendorff's $\alpha = 0.529$ (interval scale, 3 raters). Binary agreement on the sufficient/not-sufficient threshold ($\geq 4$) ranges from 76.6\% to 82.8\% across rater pairs; within-one-level agreement ranges from 81.2\% to 95.3\%. We note that the binary threshold, which drives revision-despite-sufficiency, is the better-validated metric; ordinal means should be interpreted with the moderate reliability in mind.

\subsection{Meta-Commentary Stripping}
\label{sec:methods-stripping}

Models frequently add boilerplate preambles (``Here's my revised version...'') and postambles (``Let me know if you'd like any changes'') to revision outputs. In our 50 balanced-panel pairs, 84\% of revision-side outputs contain at least one such wrapper, compared to 14\% of Turn~1 outputs. This asymmetry creates a measurement confound: LLM judges can identify the revision side by its meta-wrapping, and pointwise evaluators may trigger the ``Overdone'' criterion on boilerplate rather than content.

To control for this, we strip meta-commentary from all 3,600 outputs using validated regex patterns matching common preamble and postamble phrases. Of the 3,600 outputs, 1,200 (33.3\%) were modified by stripping and rescored by Claude Sonnet~4 at temperature~0 using the identical evaluation prompt. Validation: on the 50 balanced-panel pairs, two human annotators rating stripped content agreed with the LLM judge at $\kappa = 0.569$, compared to near-chance agreement ($\kappa$ between $-0.07$ and $+0.02$) on unstripped content. All quality results in Section~\ref{sec:results} report stripped scores as primary estimates.

\subsection{Metrics}

\textbf{Revision yield} is the mean evaluator level at each turn, restricted to genuine
revisions. \textbf{Revision direction} records, for each genuine revision, whether its level
is lower than, equal to, or higher than the level of the most recent turn whose content was
genuinely new; the comparison is made against the last genuinely new content because the
immediately preceding turn is frequently a meta-response restating the output it
accompanies. \textbf{Revision-despite-sufficiency} is the fraction of turns rated
$\geq 4$ where the model revised at the next turn. The \textbf{revision tax} is the extra
tokens spent past the quality-optimal turn $t^*$ as a percentage of optimal tokens,
$(T_{\text{full}} - T_{t^*}) / T_{t^*} \times 100$.

\subsection{Sub-experiments}

\paragraph{Targeted Feedback.} For outputs rated level~1--3, a second model instance receives the output plus a specific critique and produces a targeted revision, which is then blind-evaluated. This tests whether specific feedback outperforms the generic revision probe ($n = 177$ paired comparisons after filtering to cases where the generic next-turn output is a genuine revision, not a meta-response).

\paragraph{Reversibility.} Two human annotators independently compared stripped Turn~1 and
last-genuine-revision outputs for the 50 balanced-panel trials, blind to turn identity and
position-randomised, with a success bar of $\geq$65\% Turn~1 preference set in advance
($\kappa = 0.703$, $n = 50$). The outcome is reported in the Limitations.

